% !TeX root = ../papers/nc-p3.tex
\section{階数15のPauli基底と一般窓の有限同型判定}
\label{sec:rank-fifteen-pauli-comparison}
\begin{theorem}\label{thm:rank-fifteen-pauli-comparison}
第一表示の通常な階数15窓で、節\ref{sec:rank-fifteen-endpoint-orbits}の
$H\in M_4\otimes M_4$ は、両行列空間のPauli基底を対応させる十五項の和になる。
自然な十六次元空間を $H$ と単位元から復元できる非空開集合上では、
二窓の添字保持同型は48個のラベル置換について積の係数を比較すれば決定できる。
一般窓の独立共役に関する射影安定化群は位数16である。
\end{theorem}
\begin{proof}
$r^2=z,d^2=\lambda$ とし、$X_j$ を四文字への $j$ の平行移動とする。
$D(v_0,v_1,v_2,v_3)$ は対角行列を表す。出力と入力の順に
\[
 \begin{array}{c|cc}
     & A & B\\\hline
 0&D(r,r,1,z)X_1&D(1,1/r,d,rd)X_2\\
 1&D(1/r,r,1,1)X_1&D(1,r,d,rd)X_3
 \end{array}
\]
を取る。全四位置の $X_1$ は
\path{computations/verify_rank_fifteen_word_symmetries.py} の対角係数を伴って元のテンソルを保つ。
$B$ は定理\ref{thm:rank-fifteen-word-symmetries}の交互平行移動である。
実際 $a+e=2$ と階数低下式から $ae+zbf=2$ であり、交換式から
$q^2=u(zb^2-a^2)=-1/(\lambda z)$ となる。$q=i/(rd)$ を取れば表の係数を得る。
二つの位置のそれぞれで $A,B$ は可換で各平方はスカラーである。
また文字保存から任意の対角符号
$Z_b(j)=(-1)^{b\cdot j}$ の両側同時作用も使える。ここで $a,b,j\in\mathbb F_2^2$、
整数表示では下位ビットを第一座標とする。

各行列空間で $A,B$ の積に $Z_b$ を掛けると、物理的な文字移動 $a$ と位相 $b$ で指定される
十六本の基底直線を得る。同じ $a$ では四つの文字指標が独立であり、異なる $a$ の
台は交わらない。$A,B,Z$ の共役作用はこの空間を相異なる十六指標に分ける。
それぞれの共役の組は $H$ をスカラー倍まで保つ。入力のスカラー行列が $T$ の全核で、
出力のトレース零空間が全像であるから、そのスカラー指標は自明である。
実際、指標が非自明なら出力側の自明指標に対応する非自明な入力直線も核に入る。

$\pi$ を $1\mapsto1,\ 2\mapsto3$ の線型置換とし、$\beta=\pi^{-\mathsf T}$ とする。
対角符号から両側の移動は同じ $a$ で、$A,B$ の指標から位相はそれぞれ $b,\beta b$ となる。
従って $H$ はこの順で対応するPauli直線のテンソル積に分解する。
ラベル $x=(a,b)\in\mathbb F_2^4$ の成分を $H_x$ とし、便宜上
$H_0=I\otimes I$ と置く。$H$ 自身は非零ラベルの和で、十五成分は階数15により全て非零である。

二つの積を通常の積と $(A\otimes B)\circ(C\otimes D)=AC\otimes DB$ とする。
十六本の空間はどちらでも閉じ、各積のラベルは和になる。
$I,H$ を含み両積で閉じる最小の線型空間が十六次元なら、それがこの空間を内在的に決定する。
その可分テンソルの直線は基底の十六本だけである。実際、両Pauli基底で係数行列は対角である。
従って添字保持同型はこれらを置換し、通常積の保存から置換 $F$ は線型である。
左右の行列因子での交換子の符号も保つ。ラベル $(a,b),(c,e)$ の符号指数は
それぞれ $a\cdot e+c\cdot b$ と $a\cdot\beta e+c\cdot\beta b$ であり、
両方を保つ $F$ は48個である。コードで四基底ベクトルの像を列挙する。

$H_xH_y=c_{xy}H_{x+y}$、$e_x=\boldsymbol1_{x\ne0}$ とする。
二窓の同型の必要十分条件は、この48置換のいずれかとスカラー $h\ne0$ に対して
\[
 c'_{Fx,Fy}=h^{e_x+e_y-e_{x+y}}c_{xy}
\]
が全ラベルで成立することである。相異なる非零の一組から $h$ は一意に決まる。
十分性のため、各行列因子で四つのPauli生成元の平方を1に規格化し、移行先の
四生成元にも同じ規格化をする。四平方と六つの交換子が同じであり、
順序付き単項式は双方で $M_4$ の基底になる。生成元の対応は行列代数の同型を与え、内部共役で実現する。
二つの共役を適用した後で残る倍率は、上式により全ラベル群の指標である。
第一因子の交換子形式は非退化なので、一つのPauli元の共役がこの指標を実現する。
これで $H,H'$ のスカラーと独立共役による対応が得られ、節\ref{sec:rank-fifteen-adjoint-invariants}
の同値条件が適用される。

\path{computations/verify_rank_fifteen_pauli_comparison.py} は $101$ 上の
$(z,\lambda,n)=(4,9,46)$ で両積を閉じると十六次元になり、自己比較に成功する $F$ は
恒等写像だけになることを確かめる。単根なのでこの非消滅検査は標数零に持ち上がり、
不可約曲面の非空開集合で成り立つ。ラベルを固定する安定化元は各因子で指標を与え、
十五項を共通スカラー倍するには二指標の積が自明でなければならない。
一方の指標で他方が決まり、非退化性とPauli基底の生成性からちょうど十六組である。
\end{proof}
\begin{corollary}\label{cor:rank-fifteen-interior-101}
剰余標数101の付値を持つ標数零の体で、倍率が $(4,9)$ に還元されるとする。
この倍率の通常な階数15窓は両表示とも全て標準Koszul左右AS正則である。
\end{corollary}
\begin{proof}
上記コードは $P(n)$ の剰余根が $46,47,54,55$ の四単根で、全枠座標が非零であること、
全ての根で四中心商が全四始点の長さ13まで所定の次元を持つことを確かめる。
モニックな $P$ の定数項は単元なので全ての根は整かつ単元であり、
剰余はこの四つに限られる。経路・中心の次元を同時に確認したので、
命題\ref{prop:fixed-residual-regularity}の中心の持上げから主張が従う。
\end{proof}
